\chapter*{Abstract}
\addcontentsline{toc}{chapter}{Abstract}

Crop-protection tools are mostly reactive: they work only after a farmer has noticed a symptom and photographed it. This
project presents \textbf{GAIA}, a rover-assisted Generative AI system for proactive crop health monitoring. A small rover
(ESP32-CAM, Arduino Nano, DHT11, MQ-5, MQ-6 and a servo-mounted soil-moisture probe) is only a thin data-acquisition layer. The
effort is concentrated in two independent Generative AI pipelines that share one knowledge store (PostgreSQL with pgvector)
and are delivered through a Flutter application and a Python (FastAPI) backend.

\emph{Pipeline A} is a conversational assistant. A language model turns a question into a structured query plan, SQL computes exact
statistics, a hybrid retriever (vector and full-text search fused by Reciprocal Rank Fusion) returns relevant chunks, and Gemini
writes a streamed answer whose citations are validated against the retrieved context. \emph{Pipeline B} diagnoses crops: one
vision-language call turns a photograph into strict JSON, and one research agent searches the web for care or treatment
information; the result is shown with sources and written one way into the knowledge store.

To save battery, every sensor can be switched on or off. The backend records switches, readings and rover heartbeats, so data gaps
are explained by code (``MQ5 was switched off from 06:00 to 09:20'') and not guessed by the model. A Live view animates readings arriving,
windows closing, and chunks being written, embedded and stored.

The system was verified with 96 backend and 22 application tests and by compiling both firmware sketches. On a deterministic
synthetic farm, planner accuracy, numeric recall and citation accuracy were all 100\,\% on ten questions. The backend ingested about
5,000 sensor values per second and indexed a 15-minute window in about 73\,ms. Model confidence is shown as an uncalibrated
indicator, and diagnoses are phrased as possibilities.

\noindent\textbf{Keywords:} Generative AI, Retrieval-Augmented Generation, Vision-Language Model, Agriculture, IoT, pgvector, Gemini, Flutter.
\clearpage
